\section{Pour s'entrainer encore}
\begin{UPSTIexercice}{Table de multiplication}
	\UPSTIetape Écrire un programme \texttt{table\_multiplication.c} qui demande à l'utilisateur un entier \texttt{n} et affiche la table de multiplication de \texttt{n} de 1 à 10 en utilisant une boucle \texttt{for}.
	Par exemple, si l'utilisateur entre \texttt{5}, le programme affichera :
	\begin{lstlisting}[language=bash,style=console]
5 x 1 = 5
5 x 2 = 10
5 x 3 = 15
5 x 4 = 20
5 x 5 = 25
5 x 6 = 30
5 x 7 = 35
5 x 8 = 40
5 x 9 = 45
5 x 10 = 50
\end{lstlisting}
\end{UPSTIexercice}

\begin{UPSTIcorrectionP}{Table de multiplication}
	\begin{lstlisting}[language=c]
#include <stdio.h>

int main(void) {
    int n;
    printf("Entrez un entier : ");
    scanf("%d", &n);

    for (int i = 1; i <= 10; i++) {
        printf("%d x %d = %d\n", n, i, n * i);
    }

    return 0;
}
	\end{lstlisting}
\end{UPSTIcorrectionP}

\begin{UPSTIexercice}{Somme des n premiers entiers}
	\UPSTIetape Écrire un programme \texttt{somme\_n\_entiers.c} qui demande à l'utilisateur un entier \texttt{n} et calcule la somme des \texttt{n} premiers entiers (de 1 à \texttt{n}) en utilisant une boucle \texttt{for}. Le programme doit ensuite afficher le résultat.
	Par exemple, si l'utilisateur entre \texttt{5}, le programme affichera :
	\begin{lstlisting}[language=bash,style=console]
1 + 2 + 3 + 4 + 5 = 15
\end{lstlisting}
	\textbf{Utiliser deux boucles différentes pour le calcul de la somme et pour l'affichage du résultat.}
\end{UPSTIexercice}

\begin{UPSTIcorrectionP}{Somme des n premiers entiers}
	\begin{lstlisting}[language=c]
#include <stdio.h>

int main(void) {
    int n;
    printf("Entrez un entier : ");
    scanf("%d", &n);

    int somme = 0;
    for (int i = 1; i <= n; i++) {
        somme += i;
    }

    printf("1");
    for (int i = 2; i <= n; i++) {
        printf(" + %d", i);
    }
    printf(" = %d\n", somme);

    return 0;
}
	\end{lstlisting}
\end{UPSTIcorrectionP}

\begin{UPSTIexercice}{Factorielle}
	\UPSTIetape Écrire un programme \texttt{factorielle.c} qui demande à l’utilisateur un entier \texttt{n} et calcule \texttt{n!} (la factorielle de \texttt{n}) en utilisant une boucle \texttt{for}.
	Par exemple, pour \texttt{n = 5}, le programme affichera :
	\begin{lstlisting}[language=bash,style=console]
5! = 1 x 2 x 3 x 4 x 5 = 120
\end{lstlisting}
\end{UPSTIexercice}

\begin{UPSTIcorrectionP}{Factorielle}
	\begin{lstlisting}[language=c]
#include <stdio.h>

int main(void) {
    int n;
    printf("Entrez un entier : ");
    scanf("%d", &n);

    long long factorielle = 1;
    for (int i = 1; i <= n; i++) {
        factorielle *= i;
    }

    printf("%d! = 1", n);
    for (int i = 2; i <= n; i++) {
        printf(" x %d", i);
    }
    printf(" = %lld\n", factorielle);

    return 0;
}
	\end{lstlisting}
\end{UPSTIcorrectionP}

\begin{UPSTIexercice}{Carre}
	\UPSTIetape Écrire un programme \texttt{carre.c} qui demande à l'utilisateur un entier \texttt{n} et affiche un carré de \texttt{n} par \texttt{n} utilisant des astérisques (\texttt{*}). Utiliser une boucle \texttt{for} imbriquée pour générer les lignes et les colonnes du carré.
	Par exemple, si l'utilisateur entre \texttt{4}, le programme affichera :
	\begin{lstlisting}[language=bash,style=console]
****
****
****
****
\end{lstlisting}
\end{UPSTIexercice}

\begin{UPSTIcorrectionP}{Carre}
	\begin{lstlisting}[language=c]
#include <stdio.h>

int main(void) {
    int n;
    printf("Entrez un entier : ");
    scanf("%d", &n);

    for (int ligne = 0; ligne < n; ligne++) {
        for (int colonne = 0; colonne < n; colonne++) {
            printf("*");
        }
        printf("\n");
    }

    return 0;
}
	\end{lstlisting}
\end{UPSTIcorrectionP}

\begin{UPSTIexercice}{Triangle -- Niveau 1}
	\UPSTIetape Écrire un programme \texttt{triangle1.c} qui demande à l’utilisateur un entier \texttt{n} et affiche un triangle d’étoiles aligné à gauche, de hauteur \texttt{n}.
	Par exemple, pour \texttt{n = 4} :
	\begin{lstlisting}[language=bash,style=console]
*
**
***
****
\end{lstlisting}
\end{UPSTIexercice}

\begin{UPSTIcorrectionP}{Triangle -- Niveau 1}
	\begin{lstlisting}[language=c]
#include <stdio.h>

int main(void) {
    int n;
    printf("Entrez un entier : ");
    scanf("%d", &n);

    for (int ligne = 1; ligne <= n; ligne++) {
        for (int col = 0; col < ligne; col++) {
            printf("*");
        }
        printf("\n");
    }

    return 0;
}
	\end{lstlisting}
\end{UPSTIcorrectionP}

\begin{UPSTIexercice}{Triangle -- Niveau 2}
	\UPSTIetape Écrire un programme \texttt{triangle2.c} qui demande à l’utilisateur un entier \texttt{n} et affiche un triangle isocèle aligné au centre, de hauteur \texttt{n}.
	Par exemple, pour \texttt{n = 4} :
	\begin{lstlisting}[language=bash,style=console]
   *
  ***
 *****
*******
\end{lstlisting}
\end{UPSTIexercice}

\begin{UPSTIcorrectionP}{Triangle -- Niveau 2}
	\begin{lstlisting}[language=c]
#include <stdio.h>

int main(void) {
    int n;
    printf("Entrez un entier : ");
    scanf("%d", &n);

    for (int ligne = 1; ligne <= n; ligne++) {
        for (int espace = 0; espace < n - ligne; espace++) {
            printf(" ");
        }
        for (int etoile = 0; etoile < 2 * ligne - 1; etoile++) {
            printf("*");
        }
        printf("\n");
    }

    return 0;
}
	\end{lstlisting}
\end{UPSTIcorrectionP}

\begin{UPSTIexercice}{Damier}
	\UPSTIetape Écrire un programme \texttt{damier.c} qui demande deux entiers \texttt{lignes} et \texttt{colonnes} et affiche un damier en alternant \texttt{\#} et \texttt{.}.
	Par exemple, pour \texttt{lignes = 4} et \texttt{colonnes = 6} :
	\begin{lstlisting}[language=bash,style=console]
#.#.#.
.#.#.#
#.#.#.
.#.#.#
\end{lstlisting}
\end{UPSTIexercice}

\begin{UPSTIcorrectionP}{Damier}
	\begin{lstlisting}[language=c]
#include <stdio.h>

int main(void) {
    int lignes, colonnes;
    printf("Entrez le nombre de lignes : ");
    scanf("%d", &lignes);
    printf("Entrez le nombre de colonnes : ");
    scanf("%d", &colonnes);

    for (int i = 0; i < lignes; i++) {
        for (int j = 0; j < colonnes; j++) {
            if ((i + j) % 2 == 0) {
                printf("#");
            } else {
                printf(".");
            }
        }
        printf("\n");
    }

    return 0;
}
	\end{lstlisting}
\end{UPSTIcorrectionP}

\begin{UPSTIexercice}{Motif en diagonale}
	\UPSTIetape Écrire un programme \texttt{diagonale.c} qui demande un entier \texttt{n} et affiche un carré de taille \texttt{n} avec une diagonale de \texttt{*}.
	Exemple, pour \texttt{n = 5} :
	\begin{lstlisting}[language=bash,style=console]
*....
.*...
..*..
...*.
....*
\end{lstlisting}
\end{UPSTIexercice}

\begin{UPSTIcorrectionP}{Motif en diagonale}
	\begin{lstlisting}[language=c]
#include <stdio.h>

int main(void) {
    int n;
    printf("Entrez un entier : ");
    scanf("%d", &n);

    for (int i = 0; i < n; i++) {
        for (int j = 0; j < n; j++) {
            if (i == j) {
                printf("*");
            } else {
                printf(".");
            }
        }
        printf("\n");
    }

    return 0;
}
	\end{lstlisting}
\end{UPSTIcorrectionP}

\subsection*{Liste des checks disponibles}
\begin{lstlisting}[language=bash,style=console]
check50 IUT-GEII-Annecy/squelettes/branch-2025/tp3/maths/table_multiplication
check50 IUT-GEII-Annecy/squelettes/branch-2025/tp3/maths/somme_n_entiers
check50 IUT-GEII-Annecy/squelettes/branch-2025/tp3/maths/factorielle
check50 IUT-GEII-Annecy/squelettes/branch-2025/tp3/formes/triangle/niveau1
check50 IUT-GEII-Annecy/squelettes/branch-2025/tp3/formes/triangle/niveau2
check50 IUT-GEII-Annecy/squelettes/branch-2025/tp3/formes/damier
check50 IUT-GEII-Annecy/squelettes/branch-2025/tp3/formes/diagonale
\end{lstlisting}
